
\documentclass[addpoints]{exam}
\usepackage[version=4]{mhchem}
\usepackage{siunitx}
\usepackage{color}
\usepackage{amsmath}

% \printanswers

\newcommand{\event}{Energy Changes and Rates of Reaction}
\newcommand{\tournament}{}  % Replace with class name, date, or course code
\def\type{0}

\setlength\fillinlinelength{\textwidth}
\CorrectChoiceEmphasis{\color{red}\bfseries}
\SolutionEmphasis{\color{red}\bfseries}

\setlength\answerclearance{2pt}
\pagestyle{headandfoot}
\runningheadrule
\runningheader{\event}{\tournament}{\ifnum\type=1 Class Set Test \else Full Name:\hspace{1cm} \fi}
\runningfootrule
\runningfooter{}{Page \thepage\ of \numpages}{}

\begin{document}
\begin{center}
    {\huge Grade 12 Chemistry \\ \event
    \ifprintanswers
     \space Key
    \else
     \space \ifnum\type<2 Test \fi \ifnum\type=1 \textbf{(CLASS SET)} \fi \ifnum\type=2 Answer Sheet \fi
    \fi \\ \vspace{3mm}\tournament\vspace{1mm}}
\end{center}

\vspace{.25\textheight}

\begin{center}
    First Name: \ifprintanswers
    \underline{\hspace{3.5cm}}\textbf{KEY}\underline{\hspace{5.5cm}}\\\vspace{5mm}
    \else
        \underline{\hspace{10cm}}\\ \vspace{5mm}
    \fi
    Last Name: \underline{\hspace{10cm}}\\ \vspace{5mm}
\end{center}

\noindent\textbf{\underline{Directions:}}
\begin{itemize}
    \ifnum\type=1 \item \textbf{This test is a class set.} Please do not write any answers on this paper. \fi
    \item Show all work for full marks on short answer questions.
    \item Express answers to the appropriate number of significant figures with correct units.
    \item The test is designed to be completed in 75 minutes.
\end{itemize}
\begin{center}
\textit{For grading use only}\\
\multirowgradetable{1}[pages]
\end{center}

\newpage
\begin{questions}

\section*{Part A --- Multiple Choice (15 marks)}
\vspace{5mm}

\question[1] A reaction that releases heat to the surroundings has:
\begin{choices}
    \correctchoice $\Delta H < 0$
    \choice $\Delta H > 0$
    \choice $\Delta H = 0$
    \choice $\Delta H$ cannot be determined
\end{choices}

\question[1] \SI{100}{g} of water absorbs \SI{4180}{J} of heat. What is the temperature change?
($c_{\ce{H2O}} = \SI{4.18}{J\per g\per\celsius}$)
\begin{choices}
    \choice \SI{4.18}{\celsius}
    \correctchoice \SI{10.0}{\celsius}
    \choice \SI{41.8}{\celsius}
    \choice \SI{100}{\celsius}
\end{choices}

\question[1] Which equation correctly represents the standard enthalpy of formation of
liquid water?
\begin{choices}
    \choice \ce{2H2(g) + O2(g) -> 2H2O(l)}
    \correctchoice \ce{H2(g) + 1/2 O2(g) -> H2O(l)}
    \choice \ce{H2O(l) -> H2(g) + 1/2 O2(g)}
    \choice \ce{H(g) + OH(g) -> H2O(l)}
\end{choices}

\question[1] A catalyst increases the rate of a reaction by:
\begin{choices}
    \choice Increasing the temperature of the reaction
    \correctchoice Lowering the activation energy
    \choice Shifting the equilibrium toward products
    \choice Increasing the enthalpy change $\Delta H$
\end{choices}

\question[1] For the rate law $\text{rate} = k[\text{A}][\text{B}]^2$, the overall reaction
order is:
\begin{choices}
    \choice 1
    \choice 2
    \correctchoice 3
    \choice 4
\end{choices}

\question[1] What are the units of the rate constant $k$ for a second-order reaction?
\begin{choices}
    \choice $\text{s}^{-1}$
    \choice $\text{mol\,L}^{-1}$
    \correctchoice $\text{L\,mol}^{-1}\text{s}^{-1}$
    \choice $\text{L}^2\text{mol}^{-2}\text{s}^{-1}$
\end{choices}

\question[1] For a first-order reaction, which statement about the half-life is correct?
\begin{choices}
    \choice It increases as the reaction proceeds.
    \choice It equals $k/\ln 2$.
    \correctchoice It is independent of the initial concentration.
    \choice It decreases with each successive half-life.
\end{choices}

\question[1] Increasing the temperature of a reaction increases the rate primarily because:
\begin{choices}
    \choice The concentration of reactants increases.
    \choice The activation energy decreases.
    \correctchoice A greater fraction of molecules have energy equal to or greater than $E_a$.
    \choice The collision frequency decreases.
\end{choices}

\question[1] Using average bond enthalpies, if breaking all bonds in the reactants requires
\SI{800}{kJ} and forming all bonds in the products releases \SI{950}{kJ}, what is $\Delta H$?
\begin{choices}
    \choice $+1750$ kJ
    \choice $+150$ kJ
    \correctchoice $-150$ kJ
    \choice $-1750$ kJ
\end{choices}

\question[1] Hess's Law states that:
\begin{choices}
    \choice $\Delta H$ depends on the number of reaction steps.
    \choice $\Delta H$ is always negative for spontaneous reactions.
    \correctchoice $\Delta H$ for a reaction is independent of the pathway taken.
    \choice Catalysts change the value of $\Delta H$.
\end{choices}

\question[1] On a potential energy diagram, an endothermic reaction is characterized by:
\begin{choices}
    \choice Products at a lower energy level than reactants.
    \choice No activation energy barrier.
    \correctchoice Products at a higher energy level than reactants.
    \choice $\Delta H = 0$.
\end{choices}

\question[1] Which of the following does \textit{not} affect the rate of a chemical reaction?
\begin{choices}
    \choice Catalyst
    \choice Temperature
    \correctchoice Equilibrium constant ($K_{eq}$)
    \choice Concentration of reactants
\end{choices}

\question[1] The minimum kinetic energy that colliding molecules must possess for a reaction to
occur is called the:
\begin{choices}
    \choice Enthalpy of reaction
    \correctchoice Activation energy
    \choice Bond dissociation energy
    \choice Gibbs free energy
\end{choices}

\question[1] In a coffee-cup calorimetry experiment, the temperature of the solution
\textit{decreases} after mixing two solutions. The reaction is:
\begin{choices}
    \choice Exothermic
    \correctchoice Endothermic
    \choice At equilibrium
    \choice Catalyzed
\end{choices}

\question[1] According to the Arrhenius equation, as temperature increases, the rate constant
$k$:
\begin{choices}
    \choice Decreases exponentially
    \choice Remains constant
    \correctchoice Increases exponentially
    \choice Increases linearly
\end{choices}


\section*{Part B --- Short Answer (33 marks)}

\question A student dissolves \SI{4.00}{g} of \ce{NaOH} ($M = \SI{40.00}{g/mol}$) in
\SI{100.0}{g} of water in a coffee-cup calorimeter. The temperature rises from \SI{20.0}{\celsius}
to \SI{26.2}{\celsius}. Assume the specific heat capacity of the solution is
\SI{4.18}{J\per g\per\celsius} and the calorimeter absorbs negligible heat.
\begin{parts}
    \part[2] Calculate the heat released by the reaction, $q_{rxn}$.
    \part[2] Calculate the molar enthalpy of dissolution, $\Delta H_{sol}$, in kJ/mol.
    \part[2] The literature value for $\Delta H_{sol}$ of \ce{NaOH} is \SI{-44.5}{kJ/mol}.
    Calculate the percent error in the student's result.
    \part[2] Identify one source of experimental error and explain how it would affect the
    calculated $\Delta H_{sol}$.
\end{parts}
\begin{solution}[10cm]
\textbf{(a)} Total mass of solution $= 4.00 + 100.0 = \SI{104.0}{g}$.
\[
q_{soln} = mc\Delta T = (104.0)(4.18)(26.2 - 20.0) = (104.0)(4.18)(6.2) = \SI{2693}{J}
\]
Since the reaction is exothermic (temperature rises), $q_{rxn} = -\SI{2693}{J} = \mathbf{-2.693\ kJ}$.

\textbf{(b)} Moles of \ce{NaOH} $= \dfrac{4.00}{40.00} = \SI{0.100}{mol}$.
\[
\Delta H_{sol} = \frac{q_{rxn}}{n} = \frac{-2.693\ \text{kJ}}{0.100\ \text{mol}}
= \mathbf{-26.9\ kJ/mol}
\]

\textbf{(c)}
\[
\text{Percent error} = \frac{|{-26.9} - ({-44.5})|}{|-44.5|} \times 100\%
= \frac{17.6}{44.5} \times 100\% \approx \mathbf{39.6\%}
\]

\textbf{(d)} One source: heat loss to the surroundings (the calorimeter and surrounding air
absorb some heat). This would cause the measured $\Delta T$ to be smaller than the true value,
making $|q_{rxn}|$ appear smaller, and $|\Delta H_{sol}|$ would be \textit{underestimated}
(less negative than the true value). This is consistent with the large percent error observed.
\end{solution}

\question Use Hess's Law to calculate $\Delta H$ for the target reaction:
\[
\ce{N2(g) + 2O2(g) -> 2NO2(g)} \qquad \Delta H = \, ?
\]
given the following thermochemical equations:
\begin{align*}
\text{(1)}\quad & \ce{N2(g) + O2(g) -> 2NO(g)} \quad \Delta H_1 = +\SI{180.5}{kJ} \\
\text{(2)}\quad & \ce{2NO2(g) -> 2NO(g) + O2(g)} \quad \Delta H_2 = +\SI{113.0}{kJ}
\end{align*}
\begin{parts}
    \part[3] Show clearly how you manipulate and combine equations (1) and (2) to obtain the
    target reaction. State Hess's Law in one sentence.
    \part[2] Calculate $\Delta H$ for the target reaction.
    \part[3] Using standard enthalpies of formation ($\Delta H^\circ_f$[\ce{NO2(g)}] =
    \SI{+33.2}{kJ/mol}; elements in standard state = 0), verify your answer using the formula
    $\Delta H^\circ_{rxn} = \sum n\Delta H^\circ_f(\text{products}) - \sum n\Delta H^\circ_f
    (\text{reactants})$.
\end{parts}
\begin{solution}[9cm]
\textbf{(a)} Hess's Law: the enthalpy change of a reaction is independent of the pathway and
equals the sum of the enthalpy changes of any series of steps that lead from reactants to products.

Manipulation: keep equation (1) as written. Reverse equation (2):
\[
\text{(2$'$)}\quad \ce{2NO(g) + O2(g) -> 2NO2(g)} \qquad \Delta H_{2'} = -\SI{113.0}{kJ}
\]
Adding (1) + (2$'$):
\[
\ce{N2(g) + O2(g) + 2NO(g) + O2(g) -> 2NO(g) + 2NO2(g)}
\]
Cancelling \ce{2NO(g)} from both sides:
\[
\ce{N2(g) + 2O2(g) -> 2NO2(g)} \checkmark
\]

\textbf{(b)}
\[
\Delta H = \Delta H_1 + \Delta H_{2'} = (+180.5) + (-113.0) = \mathbf{+67.5\ kJ}
\]

\textbf{(c)}
\[
\Delta H^\circ_{rxn} = 2\Delta H^\circ_f[\ce{NO2}] - \bigl(\Delta H^\circ_f[\ce{N2}]
+ 2\Delta H^\circ_f[\ce{O2}]\bigr) = 2(+33.2) - (0 + 0) = \mathbf{+66.4\ kJ}
\]
(Small discrepancy from rounding in the given data; both methods confirm a \textbf{positive}
$\Delta H$ --- the reaction is endothermic.)
\end{solution}

\question The following data were collected for the reaction \ce{2A + B -> products}:

\begin{center}
\begin{tabular}{cccc}
\hline
Experiment & $[\text{A}]$ (mol/L) & $[\text{B}]$ (mol/L) & Initial rate (mol/L$\cdot$s) \\
\hline
1 & 0.100 & 0.100 & $2.00 \times 10^{-3}$ \\
2 & 0.200 & 0.100 & $4.00 \times 10^{-3}$ \\
3 & 0.100 & 0.200 & $8.00 \times 10^{-3}$ \\
\hline
\end{tabular}
\end{center}

\begin{parts}
    \part[3] Determine the order of reaction with respect to A and with respect to B.
    Show your working clearly.
    \part[2] Write the rate law for this reaction and calculate the value of $k$ (with units).
    \part[3] Predict the initial rate if $[\text{A}] = \SI{0.300}{mol/L}$ and
    $[\text{B}] = \SI{0.150}{mol/L}$.
\end{parts}
\begin{solution}[9cm]
\textbf{(a)} Order with respect to A (compare Experiments 1 and 2, $[\text{B}]$ constant):
\[
\frac{\text{rate}_2}{\text{rate}_1} = \frac{k(0.200)^m(0.100)^n}{k(0.100)^m(0.100)^n}
= \left(\frac{0.200}{0.100}\right)^m = 2^m = \frac{4.00\times10^{-3}}{2.00\times10^{-3}} = 2
\]
$\therefore m = 1$. The reaction is \textbf{first order} in A.

Order with respect to B (compare Experiments 1 and 3, $[\text{A}]$ constant):
\[
\frac{\text{rate}_3}{\text{rate}_1} = \left(\frac{0.200}{0.100}\right)^n = 2^n
= \frac{8.00\times10^{-3}}{2.00\times10^{-3}} = 4
\]
$\therefore n = 2$. The reaction is \textbf{second order} in B.

\textbf{(b)} Rate law: $\boxed{\text{rate} = k[\text{A}][\text{B}]^2}$

Using Experiment 1:
\[
k = \frac{\text{rate}}{[\text{A}][\text{B}]^2} = \frac{2.00\times10^{-3}}{(0.100)(0.100)^2}
= \frac{2.00\times10^{-3}}{1.00\times10^{-3}} = \mathbf{2.00\ L^2\ mol^{-2}\ s^{-1}}
\]

\textbf{(c)}
\[
\text{rate} = (2.00)(0.300)(0.150)^2 = (2.00)(0.300)(0.0225)
= \mathbf{1.35 \times 10^{-2}\ mol/L \cdot s}
\]
\end{solution}

\question The rate constants for a certain reaction are $k_1 = \SI{1.50e-3}{L\,mol^{-1}s^{-1}}$
at $T_1 = \SI{25}{\celsius}$ and $k_2 = \SI{8.20e-3}{L\,mol^{-1}s^{-1}}$ at
$T_2 = \SI{55}{\celsius}$.
\begin{parts}
    \part[3] Use the Arrhenius equation in the form
    \[
    \ln\frac{k_2}{k_1} = \frac{E_a}{R}\left(\frac{1}{T_1} - \frac{1}{T_2}\right)
    \]
    to calculate the activation energy $E_a$ for the reaction.
    ($R = \SI{8.314}{J\,mol^{-1}K^{-1}}$)
    \part[2] Calculate the rate constant at $T_3 = \SI{40}{\celsius}$ using your value of $E_a$.
    \part[3] Sketch a labelled potential energy diagram for this reaction, assuming
    $\Delta H = \SI{-45}{kJ/mol}$. Label the reactants, products, transition state, $E_a$
    (forward), $E_a$ (reverse), and $\Delta H$. Indicate whether the forward reaction is
    endothermic or exothermic.
\end{parts}
\begin{solution}[10cm]
\textbf{(a)} Convert temperatures to Kelvin: $T_1 = 298\ \text{K}$, $T_2 = 328\ \text{K}$.
\[
\ln\frac{k_2}{k_1} = \ln\frac{8.20\times10^{-3}}{1.50\times10^{-3}} = \ln(5.47) = 1.699
\]
\[
E_a = R \cdot \ln\frac{k_2}{k_1} \cdot \left(\frac{1}{T_1} - \frac{1}{T_2}\right)^{-1}
= (8.314)(1.699)\left(\frac{1}{298} - \frac{1}{328}\right)^{-1}
\]
\[
\frac{1}{298} - \frac{1}{328} = 3.356\times10^{-3} - 3.049\times10^{-3} = 3.07\times10^{-4}\ \text{K}^{-1}
\]
\[
E_a = \frac{(8.314)(1.699)}{3.07\times10^{-4}} = \frac{14.12}{3.07\times10^{-4}}
\approx \mathbf{4.60 \times 10^4\ J/mol = 46.0\ kJ/mol}
\]

\textbf{(b)} $T_3 = 313\ \text{K}$.
\[
\ln\frac{k_3}{k_1} = \frac{46000}{8.314}\left(\frac{1}{298} - \frac{1}{313}\right)
= 5532 \times (3.356\times10^{-3} - 3.195\times10^{-3}) = 5532 \times 1.61\times10^{-4} = 0.891
\]
\[
k_3 = k_1 \cdot e^{0.891} = (1.50\times10^{-3})(2.437) \approx \mathbf{3.66\times10^{-3}\ L\,mol^{-1}s^{-1}}
\]

\textbf{(c)} Diagram features:
\begin{itemize}
    \item Reactants at a higher energy level than products (exothermic, $\Delta H < 0$).
    \item A hump (transition state) above both.
    \item $E_a$ (forward) = energy from reactants to peak.
    \item $E_a$ (reverse) = energy from products to peak; $E_a(\text{rev}) = E_a(\text{fwd})
    - |\Delta H| = 46.0 - 45 = 1.0\ \text{kJ/mol}$.
    \item $\Delta H = -45\ \text{kJ/mol}$ shown as the gap between reactants and products.
\end{itemize}
The forward reaction is \textbf{exothermic}.
\end{solution}

\end{questions}
\end{document}
